\subsection{Economic interpretation}
\label{sec:review-economic-interpretation}
\begin{itemize}[itemsep=0pt,parsep=0pt,topsep=4pt]
\item \textbf{National concentration and ownership.} The event estimates, exit comparisons and focus-period results point towards an asymmetric duopoly. IndiGo and Air India Group carry most passengers, with IndiGo substantially ahead and a small competitive fringe (Tables~\ref{tab:o1-regression}, \ref{tab:o2-absorption} and~\ref{tab:o4-national-results}). Separate brands need not represent independent competitors. Tata's January 2022 acquisition of Air India raised ownership concentration through common control, without requiring passenger gains. The later mergers mainly changed brand concentration. The small partial-ownership adjustment leaves concentration classifications unchanged (Tables~\ref{tab:o2-mech} and~\ref{tab:review-o1-wedge}).

\item \textbf{National and route competition.} IndiGo and Air India Group also dominate scheduled flights on the top 20 and top 10 busiest airport pairs, but the median gap between the leading groups is smaller than nationally (Table~\ref{tab:o4-route-focus}). When approving the Air India--Vistara merger, CCI required minimum seating capacity to be maintained on seven domestic routes \parencite{cci2023vistara}. This illustrates why merger reviews also consider competition on individual airport pairs. National shares measure passengers, while airport-pair shares measure scheduled flights.

\item \textbf{Potential market power.} In the Cournot model, firms choose output while taking rivals' output as given. Applied to airlines, this provides a way to interpret competition through the capacity offered. The rise in HHI after both exits (Table~\ref{tab:o2-absorption}) would, under this model and unchanged demand elasticity, imply a higher average margin of price over marginal cost, measured as a share of price. Our share data indicate this potential for market power, but do not establish actual fare or profit increases.

Stackelberg instead describes a leader committing output before rivals respond. IndiGo's June 2023 order for 500 aircraft can be interpreted as a commitment to future capacity, but we have not tested whether rivals changed their plans in response. SpiceJet's use of vacated Jet Airways slots also illustrates why airport access matters for expansion (Appendix~\ref{app:results-o2}). Market-share gains could reflect cost advantages as well as weaker competition. Without fares and costs, we cannot distinguish these explanations, and the event estimates remain associations rather than causal effects.

\newpage
\item \label{sec:review-o3-interpretation}\textbf{Entry, scale and financial viability.} Only Vistara and Akasa met the meaningful and persistent national-share criteria. Vistara's merger left Akasa as the independent entrant at that scale (Table~\ref{tab:o3-summary}). An airline must pay fixed costs such as aircraft leases even when seats are empty. Carrying more passengers and filling more seats on existing flights spreads these costs over more tickets, reducing the fixed cost per passenger. However, a large share of a small route may still provide too few passengers to sustain the airline. Entry, revenue growth and occasional profits did not ensure survival (Tables~\ref{tab:results-o3-regional-summary}, \ref{tab:o3-fin} and~\ref{tab:results-o3-common-year}). Funding can sustain expansion despite losses. Accounting profit also need not cover investors' opportunity cost of capital. Different airline ages and incomplete finances limit these comparisons.

\item \textbf{Dependence and resilience.} Rivals' unused capacity relative to IndiGo's traffic fell from December 2024 to December 2025 (Table~\ref{tab:o4-capacity-buffer}), suggesting a thinner aggregate buffer. This rough estimate cannot show how many disrupted passengers could be accommodated, since capacity must match their routes and dates.
\end{itemize}
